\section{How the Model Was Built}\label{app:modelsteps}

The model of Section~\ref{sec:model} was not written down once. It was built in steps, each estimated by the
simulated method of moments and compared with the step before it on the same moments and the same simulation draws
(Table~\ref{tab:modelsteps}).

\begin{table}[!htbp]
\centering
\caption{How the model was built: each step, its model and the change in the criterion}
\label{tab:modelsteps}
\resizebox{\ifdim\width>\textwidth\textwidth\else\width\fi}{!}{%
\begin{tabular}{clcccc}
\toprule
Step & What changed & Model & Draws & $Q$ & $\Delta Q$ \\
\midrule
1 & first lattice: race in wage level and job loss only & A & R = 5 & 324.3 &  \\
1 & first lattice: H2 + H3 (selected) & E & R = 5 & 301.9 & $-17.0$ \\
2 & age slopes of job loss and of the on-the-job search friction & A & R = 5 & 207.4 & $-111.5$ \\
2 & same, model D & D & R = 5 & 188.3 & $-112.0$ \\
3 & re-employment in the year of a displacement & A & R = 5 & 171.6 & $-35.8$ \\
3 & same, model D & D & R = 5 & 152.8 & $-35.5$ \\
3a & match premium of $G$ jobs fixed at 0 & D & R = 5 & 155.6 & $+2.8$ \\
3b & depreciation out of work fixed at 0 (and $\mu_G = 0$) & A & R = 5 & 177.8 & $+6.2$ \\
4 & race-neutral pay-off of parental income in accumulation & A & R = 5 & 178.1 & $+0.3$ \\
4 & its racial difference (H4) & A & R = 5 & 135.3 & $-42.8$ \\
4 & H4 given H2 + H3 & E & R = 5 & 124.4 & $-34.5$ \\
5 & final lattice, H4 alone & A & R = 20 & 138.4 & $-43.6$ \\
5 & final lattice, selected model (H2 + H3 + H4) & E & R = 20 & 127.6 & $-33.0$ \\
\bottomrule
\end{tabular}}
\par\smallskip
\begin{minipage}{\textwidth}\footnotesize \textit{Notes:} $Q$: the SMM criterion on the 66 moments at the step's estimate; $\Delta Q$: against the run listed in \texttt{model/output/final\_x4/steps.json} on the same simulation draws (five simulated sons per PSID son for steps 1--4, twenty for step 5, so $Q$ is comparable only within those groups). Step 1 is the first lattice of three race hypotheses; steps 2--4 change the model one at a time; step 5 is the final lattice of Section~\ref{sec:model} (estimation draws; the ten-seed means are in Table~\ref{tab:modelhyp}). Code: \texttt{model/steps\_summary.py}.\end{minipage}
\end{table}


\paragraph{Step 1: three race hypotheses.} The first model had race in the wage level and in job loss, and three
candidate race parameters: access to good jobs (H1), the pay-off of schooling in accumulation (H2) and the entry return
to schooling (H3). Its lattice selected H2 and H3 together and rejected H1, but it reproduced only half of the opening:
a White $-$ Black gap of 0.159 at 45--49 against the PSID's 0.330, a White IGE too flat after 35 and a Black IGE that rose
with age where the data's is flat.

\paragraph{Steps 2 and 3: job loss and search by age.} Its worst misfits were not in persistence but in careers:
displacement and move rates per interview pair fall with age in the PSID and rose in the model, because biennial
interviews after 1997 lengthen the interval at older ages. Age slopes in job loss and in the on-the-job search friction
lowered $Q$ by 111 (model A) and 112 (model D). Without re-employment in the year of a displacement, every displacement
produced a full year without hours and job finding out of work went to its bound; allowing it lowered $Q$ by a further 36.
Two parameters were then fixed at zero: the match premium of good jobs, which the moments do not separate from access and
accumulation ($+2.8$), and depreciation out of work, whose free estimate was negative ($+6.2$).

\paragraph{Step 4: the pay-off of parental income.} The flat Black and rising White IGE asked for a channel through
which parental income raises accumulation for White sons. A race-neutral pay-off alone changed nothing ($+0.3$); its
racial difference (H4) lowered $Q$ by 42.8, and by 34.5 given H2 and H3.

\paragraph{Step 5: the final lattice.} All eight combinations of H1--H3 were re-estimated with H4 in the model at
twenty simulated sons per PSID son, each started from the best solution of the models nested in it, and H4 was tested
at the two ends of the lattice (Section~\ref{sec:model}). A polish re-estimation of every model from its own solution
moved no $Q$ by more than 0.12 in any of the ten runs (the eight models and the two without H4).

\paragraph{Next steps.} Two steps are coded and pending estimation: a race-specific search friction out of work,
aimed at the share of zero-hours years by race, and a job-loss rate specific to good jobs, aimed at the share in good jobs
after 35, which the model loses through moves and displacement (Section~\ref{sec:model}, untargeted validation). An
access probability that rises with age was profiled and rejected for that purpose.
